\section{Removing the incoming coefficient bound for tanh}
\label{sec:l1-compactification}
These refinements keep the same number of hidden neurons and the same
readout coefficients. In particular the readout $\ell^1$ restriction remains
unchanged. Only the original incoming coefficient bound is removed.

\begin{lemma}[Boolean compactification of one tanh ridge]
\label{lem:l1-Boolean-ridge}
Let $D=d+1$, $0<\tau\le1/2$, and
$h(x)=\tanh(v\cdot\bar x)$ on $\{-1,1\}^d$, with arbitrary finite $v$.
There is $v'$ with
\[
 \|v'\|_\infty\le (2D)^{2D}\operatorname{arctanh}(1-\tau)
\]
such that $\sup_x|\tanh(v'\cdot\bar x)-h(x)|\le\tau$.
\end{lemma}
\begin{proof}
Put $M=\operatorname{arctanh}(1-\tau)>0$. At inputs with
$|h(x)|\le1-\tau$, impose $v'\cdot\bar x=\operatorname{arctanh}h(x)$.
At the remaining inputs impose $v'\cdot\bar x\ge M$ or
$v'\cdot\bar x\le-M$ according to the sign. The original vector is
feasible. In variables $z=v'/M$, this is a finite closed polyhedron
with coefficients in $\{-1,1\}$ and right-hand sides in $[-1,1]$.

Its minimum-norm feasible point lies in the span of its active rows:
every null direction admits both sufficiently small feasible
perturbations. Select $r\le D$ independent active rows $A$, with
right-hand side $b$. The point is $A^\top(AA^\top)^{-1}b$.
The integer Gram matrix has determinant at least one and entries at
most $D$ in magnitude; its adjugate entries are bounded by
$(r-1)!D^{r-1}$. Therefore every coordinate has magnitude at most
$r^2(r-1)!D^{r-1}\le(2D)^{2D}$. If there are no active rows, the point
is zero. The output values are exact at nonsaturated inputs; at a
saturated input both values lie in the same interval of length $\tau$.
\end{proof}

\begin{lemma}[Continuous compactification of one tanh ridge]
\label{lem:l1-continuous-ridge}
Suppose the input law on $[-1,1]^d$ satisfies
\[
 P_X\{|u\cdot X-t|\le r\}\le\kappa r^\alpha
 \quad(\|u\|_2=1,\ t\in\mathbb R,\ 0<r\le1),
 \qquad \kappa\ge1,\quad0<\alpha\le1.
\]
For $0<\tau\le1/2$ and any finite $v=(b,w)$ there is $v'$ with
\begin{align*}
 \mathbb E|\tanh(v'\cdot\bar X)-\tanh(v\cdot\bar X)|&\le\tau,\\
 \|v'\|_\infty&\le
 (\sqrt d+2)\,\tfrac12\log(4/\tau)
                     (4\kappa/\tau)^{1/\alpha}.
\end{align*}
\end{lemma}
\begin{proof}
Set $M=\frac12\log(4/\tau)$,
$\delta=(\tau/(4\kappa))^{1/\alpha}\le1$, and $R=M/\delta$.
The elementary bound $1-\tanh M\le2e^{-2M}=\tau/2$ will be used.
Write $r=\|w\|_2$.

If $r\le R$ and $|b|\le R\sqrt d+M$, retain $v$. If $r\le R$ but
$|b|>R\sqrt d+M$, the original preactivation has constant sign and
magnitude exceeding $M$ everywhere; replace it by the constant
$\operatorname{sign}(b)M$. This changes the output by at most $\tau/2$.
Both cases have the asserted coefficient bound.

If $r>R$, write $u=w/r$ and $t=-b/r$, so the preactivation is
$r(u\cdot x-t)$. When $|t|>\sqrt d+\delta$, its sign is constant and
its magnitude exceeds $r\delta>M$ throughout the cube. Replace it
by the constant $-\operatorname{sign}(t)M$. Otherwise
$|t|\le\sqrt d+\delta\le\sqrt d+1$; replace the preactivation by
$R(u\cdot x-t)$. This has coefficients at most $R(\sqrt d+1)$.
Outside the slab $|u\cdot x-t|\le\delta$, both preactivations have
the same sign and magnitude at least $M$, so their tanh values differ
by at most $\tau/2$. The slab has probability at most
$\kappa\delta^\alpha=\tau/4$, where the output difference is at most
two. Thus the expected difference is at most $\tau$.
\end{proof}

For a teacher $g=\sum_{j=1}^C b_jh_j$ with $\sum_j|b_j|\le1$, apply
either lemma separately to every feature. The resulting same-width
teacher $g_\tau$ satisfies $\|g-g_\tau\|_\infty\le\tau$ in the Boolean
case, or $\mathbb E|g-g_\tau|\le\tau$ in the continuous case. Since both
functions and $f_0$ lie in $[-1,1]$,
\[
 R(g_\tau)\le R(g)+4\tau.
\]
The noise floor cancels in this comparison.

\begin{corollary}[Coefficient-free shallow tanh classes]
\label{cor:l1-coefficient-free}
Use the same posterior as in Theorem~\ref{thm:l1-main}, and fix
$N\ge8C$, $n\ge2$. The comparison class consists of
width-at-most-$C$ tanh teachers with arbitrary finite incoming
coefficients and readout $\ell^1$ norm at most one. With probability at least
$1-\delta$, its final posterior mean satisfies the following bounds.
For arbitrary Boolean-input laws, put $D=d+1$; then
\[
 \begin{aligned}
 R(\bar f_S)&\le\inf_gR(g)+C_9C/N^2\\
 &\quad+\frac{C_9s}{n}
 \left[CD\{D\log(2D)+\log(e(C+1)n)\}
                   +\log(4/\delta)\right].
 \end{aligned}
\]
Under the slab assumption of Lemma~\ref{lem:l1-continuous-ridge},
\[
 \begin{aligned}
 R(\bar f_S)&\le\inf_gR(g)+C_9C/N^2\\
 &\quad+\frac{C_9s}{n}
 \left[CD\left\{\log(eD(C+1)n)
               +\frac1\alpha\log(e\kappa n)\right\}
                    +\log(4/\delta)\right].
 \end{aligned}
\]
If also $N\ge\sqrt n$, the $C/N^2$ term is absorbed into the statistical
remainder. Constants depend only on the fixed scalar prior. No regularity or
teacher-complexity parameter is supplied to the learner. Countably
many width or regularity classes can be compared simultaneously by
the confidence allocation already stated after the main theorem.
\end{corollary}
\begin{proof}
Fix a deterministic teacher within $1/n$ of the class risk infimum
and compactify with $\tau=1/n$. In the Boolean case,
\[
 \log(B_\tau+8)\le C_{10}\{D\log(2D)+\log\log(en)\}.
\]
In the continuous case,
\[
 \log(B_\tau+8)\le C_{10}\left\{\log(2D)+\log\log(en)
                                 +\alpha^{-1}\log(e\kappa n)\right\}.
\]
Substitute these bounds in Theorem~\ref{thm:l1-main} and absorb the
$4/n$ compactification error and $1/n$ near-minimizer error. The
compactified teacher is deterministic and uses the same width and
readout variation as the original one.
\end{proof}
